\section{Requisiti}

La definizione dei requisiti di KompozeR e' stata costruita seguendo un approccio orientato all'utente. In particolare, la fase di analisi e' stata articolata a partire dall'individuazione degli stakeholder principali, dalla costruzione di personas rappresentative e dalla descrizione di scenari e casi d'uso centrali. Questo consente di giustificare i requisiti di business, funzionali e non funzionali come derivazione diretta di esigenze concrete del dominio applicativo.

\subsection{Stakeholder}

L'identificazione degli stakeholder consente di chiarire chi trae valore dal sistema, chi lo utilizza direttamente e chi influenza i vincoli progettuali.

\subsubsection{Stakeholder primari}

{\footnotesize
\renewcommand{\arraystretch}{1.2}
\begin{center}
\begin{tabular}{|p{2.2cm}|p{2.4cm}|p{2.7cm}|p{2.5cm}|}
\hline
\textbf{Stakeholder} & \textbf{Descrizione} & \textbf{Obiettivi principali} & \textbf{Bisogni critici} \\
\hline
Cliente finale & Utente che desidera progettare e acquistare una scaffalatura personalizzata & Configurare rapidamente il prodotto corretto, capire prezzo e compatibilita', completare l'acquisto senza errori & Configuratore chiaro, anteprima affidabile, salvataggio configurazioni \\
\hline
Cliente registrato ricorrente & Utente che torna sul sito per modificare o riordinare una configurazione & Riutilizzare o modificare configurazioni salvate & Storico configurazioni, accesso rapido \\
\hline
Amministratore catalogo & Operatore che gestisce prodotti, componenti, disponibilita' e prezzi & Mantenere coerente il catalogo e rendere acquistabili solo combinazioni valide & Interfacce di gestione chiare, aggiornamenti rapidi, controllo consistenza dati \\
\hline
\end{tabular}
\end{center}
}

\subsubsection{Stakeholder secondari}

{\footnotesize
\renewcommand{\arraystretch}{1.2}
\begin{center}
\begin{tabular}{|p{2.6cm}|p{3.1cm}|p{3.4cm}|}
\hline
\textbf{Stakeholder} & \textbf{Ruolo nel progetto} & \textbf{Interesse principale} \\
\hline
Azienda Kompo & Committente del prodotto & Aumentare conversione, ridurre abbandono, migliorare esperienza d'acquisto \\
\hline
Team di sviluppo & Realizza e mantiene il sistema & Requisiti chiari, comportamento verificabile, architettura evolvibile \\
\hline
Servizio clienti & Supporta utenti con dubbi o problemi & Ridurre richieste ripetitive grazie a configuratore e assistenza automatizzata \\
\hline
Responsabile business/reportistica & Analizza andamento vendite e utilizzo del sistema & Ottenere insight utili su trend di vendita e utilizzo delle funzionalita' \\
\hline
\end{tabular}
\end{center}
}

\subsection{Personas}

Per guidare la definizione dei requisiti sono state individuate tre personas rappresentative.

\subsubsection{Marta, cliente orientata alla praticita'}

Marta ha 34 anni, e' architetta freelance e desidera arredare uno studio domestico con una soluzione modulare compatibile con lo spazio disponibile. Ha buone competenze digitali e si aspetta un configuratore guidato, capace di mostrare chiaramente il risultato finale e il prezzo complessivo. Le sue principali frustrazioni riguardano la complessita' dei cataloghi tecnici, i dubbi sulla compatibilita' tra componenti e il rischio di perdere il lavoro svolto durante la configurazione.

\subsubsection{Luca, cliente indeciso ma motivato all'acquisto}

Luca ha 28 anni, e' impiegato e vuole arredare il soggiorno pur non conoscendo in dettaglio il dominio dei mobili componibili. Ha competenze digitali medie e necessita di un'interfaccia semplice, di un supporto contestuale sui prezzi e della possibilita' di salvare una configurazione per riprenderla in seguito. La sua esigenza principale e' evitare errori di comprensione durante il processo di acquisto.

\subsubsection{Giulia, amministratrice del catalogo}

Giulia ha 41 anni ed e' responsabile operativa dell'e-commerce. Si occupa di aggiornare prezzi e disponibilita' dei componenti e ha bisogno di strumenti amministrativi chiari, rapidi e coerenti con il resto del sistema. Il suo obiettivo e' garantire che le configurazioni costruite dagli utenti restino allineate con i dati reali del catalogo e con le informazioni operative.

\subsection{Scenari d'uso principali}

Gli scenari d'uso permettono di rappresentare in modo narrativo le interazioni rilevanti da cui derivano i requisiti del sistema.

\begin{itemize}
	\item \textbf{Configurazione iniziale}: l'utente seleziona una categoria di prodotto, inserisce eventuali preferenze iniziali e compone la scaffalatura attraverso il configuratore visuale. Il sistema mostra soltanto componenti compatibili e aggiorna dinamicamente anteprima e prezzo.
	\item \textbf{Supporto contestuale}: durante la configurazione l'utente apre il chatbot e richiede chiarimenti sul prezzo di uno o piu' componenti, senza interrompere il flusso operativo principale.
	\item \textbf{Salvataggio e ripresa}: l'utente autenticato salva una configurazione, la recupera in una sessione successiva e la modifica prima dell'acquisto.
	\item \textbf{Notifica di variazione}: il sistema informa l'utente quando una configurazione salvata o attiva viene impattata da una variazione di prezzo o disponibilita' di uno dei componenti coinvolti.
	\item \textbf{Gestione amministrativa}: l'amministratore aggiorna dati di catalogo e consulta ordini e report sintetici per monitorare l'andamento operativo e commerciale.
\end{itemize}

\subsection{Casi d'uso principali}

I casi d'uso principali del sistema sono riportati di seguito in forma sintetica.

\subsubsection{UC1 -- Configurazione di una scaffalatura}

\begin{description}
	\item[Attore principale] Cliente finale.
	\item[Precondizioni] Il configuratore e' disponibile e il sistema dispone dei componenti necessari alla composizione.
	\item[Flusso principale] L'utente avvia il configuratore, seleziona o conferma la categoria di prodotto, inserisce eventuali vincoli iniziali, modifica la configurazione e visualizza gli aggiornamenti dinamici di anteprima e prezzo. Il sistema consente esclusivamente l'uso di componenti compatibili con la categoria scelta.
	\item[Postcondizioni] Esiste una configurazione coerente, pronta per essere salvata o aggiunta al carrello.
\end{description}

\subsubsection{UC2 -- Richiesta di assistenza contestuale}

\begin{description}
	\item[Attore principale] Cliente finale.
	\item[Precondizioni] L'utente si trova in una fase di configurazione e il chatbot e' disponibile.
	\item[Flusso principale] L'utente apre il chatbot, formula una domanda sul prezzo di uno o piu' componenti e riceve una risposta contestuale che gli consente di proseguire la configurazione.
	\item[Postcondizioni] L'utente ottiene supporto senza uscire dal configuratore.
\end{description}

\subsubsection{UC3 -- Salvataggio e ripresa di una configurazione}

\begin{description}
	\item[Attore principale] Utente autenticato.
	\item[Precondizioni] L'utente ha creato almeno una configurazione ed e' autenticato.
	\item[Flusso principale] L'utente salva la configurazione corrente, la ritrova nel proprio profilo in una sessione successiva e la riapre per modificarla o completarla.
	\item[Postcondizioni] La configurazione risulta persistita e riutilizzabile in sessioni future.
\end{description}

\subsubsection{UC4 -- Notifica di variazione prezzo o disponibilita'}

\begin{description}
	\item[Attore principale] Utente autenticato.
	\item[Attore secondario] Amministratore catalogo.
	\item[Precondizioni] Esiste almeno una configurazione attiva o salvata e uno dei componenti associati subisce una variazione rilevante.
	\item[Flusso principale] L'amministratore aggiorna il dato di catalogo, il sistema individua le configurazioni impattate e invia una notifica agli utenti coinvolti.
	\item[Postcondizioni] L'utente e' informato delle variazioni che possono influenzare la propria configurazione.
\end{description}

\subsubsection{UC5 -- Gestione amministrativa e reportistica}

\begin{description}
	\item[Attore principale] Amministratore.
	\item[Precondizioni] L'amministratore e' autenticato nell'area gestionale.
	\item[Flusso principale] L'amministratore aggiorna prezzi e disponibilita' dei componenti, consulta gli ordini e accede a una sezione di reportistica sintetica.
	\item[Postcondizioni] I dati risultano aggiornati e l'amministratore dispone di informazioni utili al monitoraggio operativo e commerciale.
\end{description}

\subsection{Requisiti di business}

\begin{description}
	\item[BR1] Il sistema deve supportare la vendita di scaffalature personalizzabili semplificando il processo di composizione rispetto alla procedura di acquisto tradizionale.
	\item[BR2] Il sistema deve aumentare il valore percepito dell'e-commerce Kompo offrendo un configuratore visuale e servizi di assistenza digitale integrati.
	\item[BR3] Il sistema deve ridurre il rischio di errori d'acquisto dovuti a scarsa comprensione di prezzi, disponibilita', compatibilita' tra componenti e stato della configurazione.
	\item[BR4] Il sistema deve consentire all'azienda di mantenere allineati configurazioni utente, dati di catalogo e informazioni operative.
	\item[BR5] Il sistema deve fornire all'area amministrativa una vista sintetica utile al monitoraggio di ordini e trend di vendita.
\end{description}

\subsection{Requisiti funzionali}

\begin{description}
	\item[FR1] Il sistema deve fornire un configuratore visuale per la composizione della scaffalatura all'interno della categoria di prodotto selezionata.
	\item[FR2] Il sistema deve aggiornare l'anteprima della scaffalatura in seguito alle modifiche effettuate nel configuratore.
	\item[FR3] Il sistema deve aggiornare il prezzo stimato della configurazione in seguito alle modifiche effettuate nel configuratore.
	\item[FR4] Il sistema deve consentire di aggiungere al carrello una configurazione completata.
	\item[FR5] Il sistema deve fornire un chatbot integrato per rispondere a richieste contestuali sui prezzi dei componenti.
	\item[FR6] Il sistema deve consentire agli utenti autenticati di salvare una configurazione associandola al proprio profilo.
	\item[FR7] Il sistema deve consentire agli utenti autenticati di visualizzare e riprendere configurazioni salvate in precedenza.
	\item[FR8] Il sistema deve notificare agli utenti autenticati variazioni di prezzo o disponibilita' che impattano configurazioni attive o salvate.
	\item[FR9] Il sistema deve impedire la composizione di scaffalature che mescolano componenti appartenenti a categorie di prodotto diverse.
	\item[FR10] Il sistema deve consentire agli amministratori di aggiornare dati di catalogo, inclusi prezzo e disponibilita' dei componenti.
	\item[FR11] Il sistema deve rendere disponibili agli amministratori una vista ordini e una sezione di reportistica sintetica sui trend di vendita.
\end{description}

\subsection{Requisiti non funzionali}

\begin{description}
	\item[NFR1] L'interfaccia del configuratore deve permettere a un utente con competenze digitali medie di comprendere le azioni principali senza formazione preventiva.
	\item[NFR2] Il sistema deve mostrare feedback immediato a seguito delle azioni rilevanti dell'utente, in particolare durante configurazione, salvataggio e aggiornamento prezzo.
	\item[NFR3] Le principali funzionalita' del sistema devono essere progettate in modo coerente con i principi base di accessibilita' web e con l'uso corretto di elementi semantici.
	\item[NFR4] L'interfaccia deve risultare fruibile sia da dispositivi desktop sia da dispositivi mobili, preservando l'accesso alle funzionalita' essenziali.
	\item[NFR5] L'aggiornamento di anteprima e prezzo nel configuratore deve essere percepito dall'utente come tempestivo durante l'interazione.
	\item[NFR6] Le notifiche di variazione prezzo o disponibilita' devono essere recapitate in modo consistente agli utenti interessati quando il contesto applicativo lo consente.
	\item[NFR7] Le funzionalita' amministrative devono essere accessibili solo ad utenti autenticati con ruolo autorizzato.
	\item[NFR8] Il sistema deve memorizzare soltanto i dati necessari alla gestione di account, configurazioni salvate e operazioni applicative previste.
	\item[NFR9] Il sistema deve mantenere separati i moduli relativi a configurazione, gestione catalogo, notifiche e reportistica in modo da facilitare evoluzione e manutenzione.
	\item[NFR10] Il sistema deve evitare funzionalita' superflue e privilegiare interazioni essenziali, riducendo complessita', traffico e carico computazionale non necessario.
\end{description}

\subsection{Tracciabilita'}

I requisiti presentati in questa sezione derivano dall'analisi degli stakeholder, dalle personas individuate e dagli scenari e casi d'uso principali. Questa struttura garantisce la coerenza tra bisogni degli utenti, comportamento atteso del sistema e qualita' richieste al prodotto finale.

